% Figure: diffraction grating. Incident plane wave on a grating of period d;
% diffracted orders m = -2..2 at angles set by d sin(theta_m) = m lambda.
% Geometry illustrative, not to scale.
\begin{figure}[htbp]
\centering
\begin{tikzpicture}[>=Stealth, scale=1.0]
% grating plane (vertical) with tick marks for the rulings
\draw[engblue, very thick] (0,-2.3) -- (0,2.3);
\foreach \yy in {-2,-1.5,-1,-0.5,0,0.5,1,1.5,2}
  {\draw[engblue, thick] (-0.12,\yy) -- (0.12,\yy);}
\node[engblue, font=\small, anchor=south] at (0,2.3) {grating, period $d$};
% incident plane wave from left (normal incidence)
\foreach \yy in {-1.5,-0.75,0,0.75,1.5} {\draw[engred, ->] (-2.4,\yy) -- (-0.3,\yy);}
\node[engred, font=\scriptsize, anchor=south] at (-1.6,1.6) {incident};
% diffracted orders from origin: m=0 straight, m=+-1, m=+-2 at increasing angle
% angles chosen for clarity: m1 = 22 deg, m2 = 48 deg
\def\Lr{4.6}
% m=0
\draw[enggray, very thick, ->] (0,0) -- (\Lr,0);
\node[enggray, font=\scriptsize, anchor=west] at (\Lr,0) {$m=0$};
% m=+1 (22 deg up)
\draw[engorange, thick, ->] (0,0) -- (22:\Lr);
\node[engorange, font=\scriptsize, anchor=south west] at (22:\Lr) {$m=+1$};
% m=-1
\draw[engorange, thick, ->] (0,0) -- (-22:\Lr);
\node[engorange, font=\scriptsize, anchor=north west] at (-22:\Lr) {$m=-1$};
% m=+2 (48 deg)
\draw[englime!70!black, thick, ->] (0,0) -- (48:\Lr);
\node[englime!70!black, font=\scriptsize, anchor=south west] at (48:\Lr) {$m=+2$};
% m=-2
\draw[englime!70!black, thick, ->] (0,0) -- (-48:\Lr);
\node[englime!70!black, font=\scriptsize, anchor=north west] at (-48:\Lr) {$m=-2$};
% angle arc for m=+1
\draw[engorange] (1.4,0) arc [start angle=0, end angle=22, radius=1.4];
\node[engorange, font=\scriptsize] at (24:1.7) {$\theta_1$};
\end{tikzpicture}
\caption[Diffraction grating orders]{A diffraction grating of period $d$ at
normal incidence sends light into discrete orders at angles set by the
grating equation $d\sin\theta_m = m\lambda$. The zeroth order ($m=0$) passes
straight through; higher orders spread to larger angles, and the angle of a
given order grows with wavelength, which is why a grating disperses white
light into a spectrum. The number of visible orders is limited by
$|m| \le d/\lambda$.}
\label{fig:vol04ch11:grating-orders}
\end{figure}
